%%%%%%%%%%%%%%%%%%%%%%%%%%%%%%%%%%%%%%%%%%%%%%%%%%%%%%%%%%%%%%%%
% General utilities
%%%%%%%%%%%%%%%%%%%%%%%%%%%%%%%%%%%%%%%%%%%%%%%%%%%%%%%%%%%%%%%%
\usepackage[none]{hyphenat}
\pretolerance=10000
\usepackage{multirow}

\makeatletter
\newcommand{\manuallabel}[2]{\def\@currentlabel{#2}\label{#1}}
\makeatother

\newcommand{\unref}[1]{Unit~\ref{#1}}

\newcounter{halfcols}
\setcounter{halfcols}{\value{cols}/2}

\usepackage{textcomp}
\newcommand{\degree}{\textdegree}

\usepackage{ifthen}
\usepackage{forloop}% http://ctan.org/pkg/forloop
\newcounter{loopcntr}
\newcommand{\rpt}[2]{%
  \forloop{loopcntr}{0}{\value{loopcntr}<#1}{#2}%
}

\usepackage{multirow}

%\usepackage{enumerate}
\usepackage{mdwlist}
\let\note\undefined
\usepackage{calc}
\usepackage{multicol}
\setlength{\columnseprule}{0pt}

\renewenvironment{itemize*}{\begin{itemize}[nosep]}{\end{itemize}}
\renewenvironment{enumerate*}{\begin{enumerate}[nosep]}{\end{enumerate}}

\newcommand{\setabc}{%
\renewcommand{\theenumi}{\alph{enumi}}%
\renewcommand{\labelenumi}{\theenumi.}%
}

\usepackage{xparse}

\newcommand{\nic}{\noindent\centering}

\newcommand{\noi}{\noindent}

\newcommand{\newstuff}[1]{{\color{red}#1}}
\newcommand{\oldstuff}[1]{{\color{blue!50!green}#1}}
%%%%%%%%%%%%%%%%%%%%%%%%%%%%%%%%%%%%%%%%%%%%%%%%%%%%%%%%%%%%%%%%
% Theorems and nifty math
%%%%%%%%%%%%%%%%%%%%%%%%%%%%%%%%%%%%%%%%%%%%%%%%%%%%%%%%%%%%%%%%
\usepackage{amsmath,amsthm,amsfonts,amssymb}
%\usepackage[framemethod=tikz]{mdframed}
\renewcommand{\emptyset}{\varnothing}

%Let's do that again - build to mimic old mdframed
%now using tcolorbox
\usepackage{tcolorbox}
\tcbset{colback=white,
				colframe=black,
				colbacktitle=white,
				coltitle=black,
				fonttitle=\bfseries,
				boxrule=0.4pt,
				sharp corners=all,
				left*=10pt,right*=10pt,
				top=2.5pt,bottom=2.5pt,
				toptitle=2.5pt,bottomtitle=2.5pt,
				adjust text=ABC}

\NewDocumentEnvironment{mdframed}{o +b}{%
	\IfValueTF{#1}{%
		\begin{tcolorbox}[adjusted title={#1}]%
	}{%
		\begin{tcolorbox}%
	}%
	#2
}{%
	\end{tcolorbox}%
	\noindent%
}

\NewDocumentEnvironment{mdtitle}{o m +b}{%
\begin{mdframed}[\IfValueTF{#1}{#2: #1}{#2}]%
#3%
\end{mdframed}%
}{}

\newenvironment{displaybox}[1][]{\begin{mdframed}[#1]}{\end{mdframed}}

\NewDocumentEnvironment{theorem}{o +b}{\begin{mdtitle}[#1]{Theorem}#2}{\end{mdtitle}}
\NewDocumentEnvironment{corollary}{o +b}{\begin{mdtitle}[#1]{Corollary}#2}{\end{mdtitle}}
\NewDocumentEnvironment{definition}{o +b}{\begin{mdtitle}[#1]{Definition}#2}{\end{mdtitle}}
\NewDocumentEnvironment{procedure}{m +b}{\begin{mdtitle}[#1]{Procedure}#2}{\end{mdtitle}}

\NewDocumentEnvironment{fact}{o +b}{\begin{mdframed}[#1] #2}{\end{mdframed}}

\newcommand{\term}[1]{%
\boldmath%
\textbf{#1}%
\unboldmath%
}

\theoremstyle{definition}
\newtheorem*{example}{Example}
\newtheorem*{note}{Note}

\newcommand{\rrt}[2]{$\pm\frac{#1}{#2}$}

\DeclareMathOperator{\rng}{Rng}
\DeclareMathOperator{\dom}{Dom}

\newenvironment{system}{%
	\displaystyle\left\{
	\begin{array}{l}
}{
	\end{array}
	\right.
}

\usepackage{array}
\newcolumntype{o}{@{}>{{}}c<{{}}@{}}

\newcounter{eqnLbl}
\newcommand{\lbleqn}{%
	\addtocounter{eqnLbl}{1}
	(\Roman{eqnLbl})
}
\newenvironment{systemL}{%
	\setcounter{eqnLbl}{0}
	\displaystyle\left\{
	\begin{array}{l@{\qquad}c<{\lbleqn}}
}{
	\end{array}
	\right.
}
\newenvironment{systemLa}{%
	\setcounter{eqnLbl}{0}
	\displaystyle\left\{
	\begin{array}{rororor@{\qquad}c<{\lbleqn}}
}{
	\end{array}
	\right.
}

\newenvironment{variables}{%
	\begin{center}\begin{tabular}{>{$}r<{{}={}$}@{}l}%
}{%
	\end{tabular}\end{center}%
}

\DeclareFontFamily{OT1}{pzc}{}
\DeclareFontShape{OT1}{pzc}{m}{it}%
{<-> s * [1.15] pzcmi7t}{}
\DeclareMathAlphabet{\mathscr}{OT1}{pzc}{m}{it}
\newcommand{\emm}{\mathscr{m}}
\newcommand{\enn}{\mathscr{n}}

\newcommand{\pw}[1]{%
\left\{
\begin{array}{r@{\qquad\mbox{if}\quad}roror}
#1
\end{array}
\right.
}

\newcommand{\bbR}{\mathbb{R}}
\newcommand{\bbN}{\mathbb{N}}
\newcommand{\bbZ}{\mathbb{Z}}
\newcommand{\bbW}{\mathbb{W}}
\newcommand{\bbQ}{\mathbb{Q}}

\newcommand{\relR}{\mathrel{R}}
\newcommand{\relS}{\mathrel{S}}
\newcommand{\relT}{\mathrel{T}}

%%%%%%%%%%%%%%%%%%%%%%%%%%%%%%%%%%%%%%%%%%%%%%%%%%%%%%%%%%%%%%%%
% Set notation etc
%%%%%%%%%%%%%%%%%%%%%%%%%%%%%%%%%%%%%%%%%%%%%%%%%%%%%%%%%%%%%%%%
\newcommand{\Set}[1]{\left\{#1\right\}}
\newcommand{\N}{\ensuremath{\mathbb{N}}}
\newcommand{\Z}{\ensuremath{\mathbb{Z}}}
\newcommand{\Q}{\ensuremath{\mathbb{Q}}}
\newcommand{\R}{\ensuremath{\mathbb{R}}}

\newcommand{\rel}[1]{\mathbin{#1}}

\newcommand{\Abs}[1]{\left|#1\right|}


%%%%%%%%%%%%%%%%%%%%%%%%%%%%%%%%%%%%%%%%%%%%%%%%%%%%%%%%%%%%%%%%
% Picture environments
%%%%%%%%%%%%%%%%%%%%%%%%%%%%%%%%%%%%%%%%%%%%%%%%%%%%%%%%%%%%%%%%
\newlength{\psize}

\newenvironment{colparts}[2][]{%
\begin{center}
\begin{tikzpicture}[#1,x=\textwidth/#2-1pt/#2]
}{%
\end{tikzpicture}
\end{center}
}

\newenvironment{cpic}[3][]{%
\begin{center}
\begin{tikzpicture}[#1]
\draw[step=\linespace,npblue,use as bounding box]
    (-\value{cols}\linespace,0) grid (\value{cols}\linespace,-#3);
}{%
\end{tikzpicture}
\end{center}
}

\newenvironment{lpic}[3][]{%
\begin{center}
\begin{tikzpicture}[#1]
\draw[step=\linespace,npblue,use as bounding box]
    (0,0) grid (2*\value{cols}\linespace,-#3);
}{%
\end{tikzpicture}
\end{center}
}

\newenvironment{lpicsplit}[3][]{%
\begin{center}
\begin{tikzpicture}[#1]
\useasboundingbox (0,0) rectangle (2*\value{cols}\linespace+\linespace,-#3);
\draw[step=\linespace,npblue]
    (0,0) grid ++(\value{cols}\linespace,-#3);
\draw[step=\linespace,npblue]
    (\value{cols}\linespace+\linespace,0) grid ++(\value{cols}\linespace,-#3);
}{%
\end{tikzpicture}
\end{center}
}

\newcounter{fcols}
\newenvironment{lpicvsplit}[3][]{%
\setcounter{fcols}{1 * \ratio{\paperwidth - \lrmargin - \rmargin - #2\linespace + \linespace}{#2\linespace}}%
\begin{center}
\begin{tikzpicture}[#1]
\useasboundingbox (0,0) rectangle (#2*\value{fcols}*\linespace+#2*\linespace-\linespace,-#3);
\foreach \x in {1,...,#2} {
	\draw[step=\linespace,npblue]
			(\x*\value{fcols}+\x-\value{fcols}-1,0) grid ++(\value{fcols}\linespace,-#3);
}
}{%
\end{tikzpicture}
\end{center}
}

\newcommand{\regtextcolc}[3]{\regtext{#1*\value{fcols}+#1+0.5*\value{fcols}}{#2}{#3}}
\newcommand{\regtextcoll}[3]{\regtextleft{#1*\value{fcols}+#1}{#2}{#3}}
\newcommand{\regtextcolr}[3]{\regtextright{#1*\value{fcols}+#1+\value{fcols}}{#2}{#3}}

\newcommand{\why}[2][\value{cols}]{%
	\begin{lpic}{}{#2}
	\draw (0,-1) -- (2*\value{cols},-1);
	\draw (#1,0) -- (#1,-#2);
	\foreach \x/\lbl in {0/Step,#1/Why?}
		\regtextleft{\x+0.5}{-1}{\lbl};
	\end{lpic}
}

\newlength{\aboveline}
\newlength{\fracline}
\setlength{\aboveline}{3pt}
\setlength{\fracline}{-3pt}

\edef\measurepage{\the\dimexpr\pagegoal-\pagetotal-\baselineskip\relax}

\newcounter{rowsleft}
\newlength{\spaceleft}

\newcommand{\fpcount}{%
	\setlength{\spaceleft}{\textheight-\pagetotal-\baselineskip}
	\setcounter{rowsleft}{1*\ratio{\spaceleft}{\linespace}}
}
\newcommand{\fillpaper}{%
	\fpcount
	\begin{lpic}{}{\value{rowsleft}}\end{lpic}
}
\newenvironment{filler}{%
	\fpcount%
	\begin{lpic}{}{\value{rowsleft}}%
}{%
	\end{lpic}%
}


%%%%%%%%%%%%%%%%%%%%%%%%%%%%%%%%%%%%%%%%%%%%%%%%%%%%%%%%%%%%%%%%
% Guided notes commands
%%%%%%%%%%%%%%%%%%%%%%%%%%%%%%%%%%%%%%%%%%%%%%%%%%%%%%%%%%%%%%%%
%\newcommand{\makeblank}[1][2in]{\rule{#1}{0.4pt}}
\newcommand{\makeblank}[1][2in]{%
\begin{tikzpicture}[baseline={([yshift=\aboveline]current bounding box.south)}]
\draw[npblue] (0,0) -- (#1,0);
\end{tikzpicture}
}

\newcommand{\htstrut}[1][2ex]{\rule{0px}{#1}}

\newcommand{\lbr}[1][1]{\\[#1ex]}

\newcommand{\makeblanksmall}{\makeblank[.5in]}
\newcommand{\makeblanknew}{\lbr\makeblank[5in]}
\newcommand{\makeblanknewf}{\lbr\makeblank[\linewidth-1pt]}
\newcommand{\makeblankex}{\lbr\textbf{Example:\ }\makeblank[4in]}
\newcommand{\notebreak}{\lbr\textbf{Note:\ }}

\newlength{\blanktemp}
\newcommand{\makeblankfill}[2][\linewidth] {
\setlength{\blanktemp}{#1-\widthof{#2 }-1pt}
#2 \makeblank[\blanktemp]
}


\newcommand{\blankpair}{\ensuremath{\bigl(\makeblank[.25in],\makeblank[.25in]\bigr)}}
\newcommand{\blanktrip}{\ensuremath{\bigl(\makeblank[.25in],\makeblank[.25in],\makeblank[.25in]\bigr)}}
\newcommand{\blankparens}[1][0.5in]{\ensuremath{\bigl(\makeblank[#1]\bigr)}}

\newcommand{\makespace}[1][1in]{\rule{#1}{0px}}
\newcommand{\figdash}{\makespace[\widthof{$-$}]}

\newcommand{\makehspace}[1][1in]{\rule{0px}{#1}}

%%%%%%%%%%%%%%%%%%%%%%%%%%%%%%%%%%%%%%%%%%%%%%%%%%%%%%%%%%%%%%%%
% Commands for creating picture environments
%%%%%%%%%%%%%%%%%%%%%%%%%%%%%%%%%%%%%%%%%%%%%%%%%%%%%%%%%%%%%%%%
\newenvironment{ilpic}[1][]{%
	\begin{tikzpicture}[baseline={([yshift=\aboveline]current bounding box.north)},
	                    #1]
	}{\end{tikzpicture}}

\newenvironment{ilgraph}{\begin{ilpic}[x=0.5\linespace,y=0.5\linespace]}{\end{ilpic}}

\newenvironment{icpic}[1][]{\begin{center}\begin{tikzpicture}[#1]}{\end{tikzpicture}\end{center}}

%%%%%%%%%%%%%%%%%%%%%%%%%%%%%%%%%%%%%%%%%%%%%%%%%%%%%%%%%%%%%%%%
% Commands for making pictures of math
%%%%%%%%%%%%%%%%%%%%%%%%%%%%%%%%%%%%%%%%%%%%%%%%%%%%%%%%%%%%%%%%
\newlength{\argwidth}

\newcommand{\regtext}[3]{%
\draw [yshift=\aboveline]%
			(#1,#2) node[anchor=base]{#3}%
}

\newcommand{\regtextml}[5][\textwidth]{%
\draw [yshift=\aboveline+#4\baselineskip]%
			(#2,#3) node[anchor=base,align=center,text width=#1]{#5}%
}

\newcommand{\bigtext}[3]{\regtext{#1}{#2}{\Large\bfseries #3}}

\newcommand{\regtextleft}[3]{%
\setlength{\argwidth}{\widthof{#3}}
\draw [yshift=\aboveline,xshift=.5\argwidth+\aboveline]%
			(#1,#2) node[anchor=base]{#3}%
}

\newcommand{\regtextleftblock}[4]{%
\draw[yshift=\aboveline,xshift=\aboveline] (#1,#2) node[anchor=base,text width=#4,align=left]{#3}%
}

\newcommand{\regtextright}[3]{%
\setlength{\argwidth}{\widthof{#3}}
\draw [yshift=\aboveline,xshift=-.5\argwidth-\aboveline]%
			(#1,#2) node[anchor=base]{#3}%
}

\newcommand{\fractext}[3]{%
\draw [yshift=\fracline]
      (#1,#2) node[anchor=base]{#3}%
}%

\newcommand{\fractextleft}[3]{%
\setlength{\argwidth}{\widthof{#3}}
\draw [yshift=\fracline,xshift=.5\argwidth+\aboveline]%
			(#1,#2) node[anchor=base]{#3}%
}

\newcommand{\fractextright}[3]{%
\setlength{\argwidth}{\widthof{#3}}
\draw [yshift=\fracline,xshift=-.5\argwidth-\aboveline]%
			(#1,#2) node[anchor=base]{#3}%
}

\newcommand{\regtextlp}[3]{\regtext{#1-\value{cols}}{#2}{#3}}
\newcommand{\regtextleftlp}[3]{\regtextleft{#1-\value{cols}}{#2}{#3}}
\newcommand{\regtextrightlp}[3]{\regtextright{#1-\value{cols}}{#2}{#3}}

\newcommand{\writethink}[3][\value{cols}]{%
	\twocol[#1]{#2}{#3}{Write:}{Think:}
}

\newcommand{\twocol}[5][\value{cols}]{%
	\twocollbl[#1]{#2}{#4}{#5}
	\draw[dashed](#1,#2+1) -- (#1,#2-#3+1);
}
\newcommand{\twocolleft}[5][\value{cols}]{%
	\twocollblleft[#1]{#2}{#4}{#5}
	\draw[dashed](#1,#2+1) -- (#1,#2-#3+1);
}

\newcommand{\twocollbl}[4][\value{cols}]{%
  \regtext{.5*#1}{#2}{#3};
	\regtext{.5*#1+\value{cols}}{#2}{#4};
}
\newcommand{\twocollblleft}[4][\value{cols}]{%
  \regtextleft{1}{#2}{#3};
	\regtextleft{#1+1}{#2}{#4};
}

\newcommand{\checkbox}[1]{%
\setlength{\argwidth}{\widthof{ #1}}
\begin{tikzpicture}[baseline={([yshift=\aboveline]current bounding box.south)}]
\useasboundingbox(-1,0) rectangle (\argwidth,0.75);
\draw[npblue] (-0.5,.1) rectangle (0,.6);
\draw (0.5\argwidth,\aboveline) node[anchor=base]{ #1};
\end{tikzpicture}
}

%%%%%%%%%%%%%%%%%%%%%%%%%%%%%%%%%%%%%%%%%%%%%%%%%%%%%%%%%%%%%%%%
% Graphing gear
%%%%%%%%%%%%%%%%%%%%%%%%%%%%%%%%%%%%%%%%%%%%%%%%%%%%%%%%%%%%%%%%
\newlength{\mxshift}
\newlength{\graphwidth}
\newlength{\graphht}
\newcounter{graphsize}
\newcounter{grleft}
\newcounter{grright}
\newcounter{grtop}
\newcounter{grbot}
\newcounter{zback}
\newcounter{zfront}
\newcounter{countbyx}
\newcounter{nextstepx}
\newcounter{countspacex}
\newcounter{countbyy}
\newcounter{nextstepy}
\newcounter{countspacey}
\newcounter{countbyz}
\newcounter{nextstepz}
\newcounter{countspacez}
\newcounter{divspacex}
\newcounter{divspacey}
\newcounter{divspacez}
\newcommand{\graphsquare}{%
	\setgraph{\value{graphsize}}
	         {\value{graphsize}}
					 {\value{graphsize}}
					 {\value{graphsize}}
	\setz{\value{graphsize}}{\value{graphsize}}
}
\newcommand{\setgraphsmall}{%
	\setcounter{graphsize}{\value{halfcols}-1}
	\graphsquare
}
\newcommand{\setgraphlarge}{%
	\setcounter{graphsize}{\value{cols}-1}
	\graphsquare
}
\newcommand{\setgraphsquare}[1]{%
	\setcounter{graphsize}{#1}
	\graphsquare
}
\newcommand{\setgraph}[4]{%
	\setcounter{grleft}{#1}
	\setcounter{grright}{#2}
	\setcounter{grbot}{#3}
	\setcounter{grtop}{#4}
	\setspace{2}
	\setlength{\graphwidth}{\value{grleft}\linespace+\value{grright}\linespace+\linespace}
	\setlength{\graphht}{\value{grtop}\linespace+\value{grbot}\linespace+\linespace}
}
\NewDocumentCommand{\setspace}{O{1} m O{1}}{%
	\setcounter{countbyx}{#2}
	\setcounter{countspacex}{#1}
	\setcounter{divspacex}{#3}
	\setcounter{countbyy}{#2}
	\setcounter{countspacey}{#1}
	\setcounter{divspacey}{#3}
	\setcounter{countbyz}{#2}
	\setcounter{countspacez}{#1}
	\setcounter{divspacez}{#3}
}
\newcommand{\setgraphlow}[1]{%
	\setcounter{grleft}{\value{cols}-1}
	\setcounter{grright}{\value{cols}-1}
	\setcounter{grbot}{#1}
	\setcounter{grtop}{2*\value{cols}-#1-2}
}
\newcommand{\setgraphlowright}[2]{%
	\setcounter{grleft}{2*\value{cols}-#1-2}
	\setcounter{grright}{#1}
	\setcounter{grbot}{#2}
	\setcounter{grtop}{2*\value{cols}-#2-2}
}
\newcommand{\setz}[2]{%
	\setcounter{zback}{#1}
	\setcounter{zfront}{#1}
}
\newcommand{\setcountspace}[1]{%
	\setcounter{countspacex}{#1}%
	\setcounter{countspacey}{#1}%
	\setcounter{countspacez}{#1}%
}

\newcommand{\graphscaley}[1]{%
	\setcounter{divspacey}{#1*\value{divspacey}}
}

\NewDocumentCommand{\setcountby}{m m O{1}}{%
	\setcounter{countbyx}{#1}%
	\setcounter{countbyy}{#2}%
	\setcounter{countbyz}{#3}%
}

\newcommand{\AxesGen}[2]{%
	\setcounter{nextstepx}{\value{countbyx}*2}
	\setcounter{nextstepy}{\value{countbyy}*2}
	%%%%%%%%%%%%%%%%%
	\draw[axis] ({(-\value{grleft}-.5*\value{divspacex})*\value{countspacex}/\value{divspacex}
								-\value{countspacex}/\value{divspacex}+1/\value{divspacex}},0) -- 
	            ({(\value{grright}+.5*\value{divspacex})*\value{countspacex}/\value{divspacex}
								+\value{countspacex}/\value{divspacex}-1/\value{divspacex}},0) 
							node[anchor=south,font=#1]{$x$};
	\draw[axis] (0,{(-\value{grbot}-.5*\value{divspacey})*\value{countspacey}/\value{divspacey}
								-\value{countspacey}/\value{divspacey}+1/\value{divspacey}}) -- 
	            (0, {(\value{grtop}+.5*\value{divspacey})*\value{countspacey}/\value{divspacey}
								+\value{countspacey}/\value{divspacey}-1/\value{divspacey}}) 
							node[anchor=west,font=#1]{$y$};
	%%%%%%%%%%%%%%%%%
	\setlength{\mxshift}{\widthof{#2$-$}}
	%%%%%%%%%%%%%%%%%
	\ifthenelse{\value{grright}<\value{nextstepx}}{
		\ifthenelse{\value{grright}<\value{countbyx}}{}{%If it's too small, do nothing
			\draw[tick] (\value{countbyx}*\value{countspacex}/\value{divspacex},.1) 
							-- (\value{countbyx}*\value{countspacex}/\value{divspacex},-.1) 
			            node[anchor=north,font=#2]{$\arabic{countbyx}$};			
		}
	}{
		\foreach \x in {\arabic{countbyx},\arabic{nextstepx},...,\value{grright}}
			\draw[tick] (\x*\value{countspacex}/\value{divspacex},.1) 
							-- (\x*\value{countspacex}/\value{divspacex},-.1)
			            node[anchor=north,font=#2]{$\x$};
	}
	%%%%%%%%%%%%%%%%%
	\ifthenelse{\value{grleft}<\value{nextstepx}}{
		\ifthenelse{\value{grleft}<\value{countbyx}}{}{%If it's too small, do nothing
			\draw[tick] (-\value{countbyx}*\value{countspacex}/\value{divspacex},.1) 
							-- (-\value{countbyx}*\value{countspacex}/\value{divspacex},-.1) 
			            node[anchor=north,font=#2]{$-\arabic{countbyx}\phm$};			
		}
	}{
		\foreach \x in {\arabic{countbyx},\arabic{nextstepx},...,\value{grleft}}
			\draw[tick] (-\x*\value{countspacex}/\value{divspacex},.1) 
							-- (-\x*\value{countspacex}/\value{divspacex},-.1)
			            node[anchor=north,font=#2]{$-\x\phm$};
	}
	%%%%%%%%%%%%%%%%%
	\ifthenelse{\value{grtop}<\value{nextstepy}}{
		\ifthenelse{\value{grtop}<\value{countbyy}}{}{%If it's too small, do nothing
			\draw[tick] (.1,\value{countbyy}*\value{countspacey}/\value{divspacey}) 
							-- (-.1,\value{countbyy}*\value{countspacey}/\value{divspacey}) 
			            node[anchor=east,font=#2]{$\arabic{countbyy}$};			
		}
	}{
		\foreach \x in {\arabic{countbyy},\arabic{nextstepy},...,\value{grtop}}
			\draw[tick] (.1,\x*\value{countspacey}/\value{divspacey}) 
							-- (-.1,\x*\value{countspacey}/\value{divspacey}) 
			            node[anchor=east,font=#2]{$\x$};
	}
	%%%%%%%%%%%%%%%%%
	\ifthenelse{\value{grbot}<\value{nextstepy}}{
		\ifthenelse{\value{grbot}<\value{countbyy}}{}{%If it's too small, do nothing
			\draw[tick] (.1,-\value{countbyy}*\value{countspacey}/\value{divspacey}) 
						-- (-.1,-\value{countbyy}*\value{countspacey}/\value{divspacey}) 
			            node[anchor=east,font=#2]{$-\arabic{countbyy}$};			
		}
	}{
		\foreach \x in {\arabic{countbyy},\arabic{nextstepy},...,\value{grbot}}
			\draw[tick] (.1,-\x*\value{countspacey}/\value{divspacey}) 
						-- (-.1,-\x*\value{countspacey}/\value{divspacey})
			            node[anchor=east,font=#2]{$-\x$};
	}
}

\newcommand{\zAxis}[2]{%
	\setcounter{nextstepz}{\value{countbyz}*2}
	\draw[axis] (-135:{(-\value{zback}-.5*\value{divspacez})*\value{countspacez}/\value{divspacez}
									-\value{countspacez}/\value{divspacez}+1/\value{divspacez}}) 
						-- (-135:{(\value{zfront}+.5*\value{divspacez})*\value{countspacez}/\value{divspacez}
									+\value{countspacez}/\value{divspacez}-1/\value{divspacez}}) 
							node[anchor=south east,font=#1]{$z$};
	%%%%%%%%%%%%%%%%%
	\ifthenelse{\value{zfront}<\value{nextstepz}}{
		\ifthenelse{\value{zfront}<\value{countbyz}}{}{%If it's too small, do nothing
			\draw[tick] (-135:\value{countbyz}*\value{countspacez}/\value{divspacez}) ++ (135:.1) 
						-- ++(135:-.2) 
			            node[anchor=north west,font=#2]{$\arabic{countbyz}$};			
		}
	}{
		\foreach \x in {\arabic{countbyz},\arabic{nextstepz},...,\value{zfront}}
			\draw[tick] (-135:\x*\value{countspacez}/\value{divspacez}) ++ (135:.1) -- ++(135:-.2) 
			            node[anchor=north west,font=#2]{$\x$};
	}
	%%%%%%%%%%%%%%%%%
	\ifthenelse{\value{zback}<\value{nextstepz}}{
		\ifthenelse{\value{zback}<\value{countbyz}}{}{%If it's too small, do nothing
			\draw[tick] (-135:-\value{countbyz}*\value{countspacez}/\value{divspacez}) ++ (135:.1) 
						-- ++(135:-.2) 
			            node[anchor=north west,font=#2]{$-\arabic{countbyz}$};			
		}
	}{
		\foreach \x in {\arabic{countbyz},\arabic{nextstepz},...,\value{zback}}
			\draw[tick] (-135:-\x*\value{countspacez}/\value{divspacez}) ++ (135:.1) -- ++(135:-.2) 
			            node[anchor=north west,font=#2]{$-\x$};
	}
}

\newcommand{\Axes}[2]{%
	\setspace{2}
	\AxesGen{#1}{#2}
	%\draw[axis] (-\value{grleft}-.5,0) -- 
	                      %( \value{grright}+.5,0) node[anchor=south,font=#1]{$x$};
	%\draw[axis] (0,-\value{grbot}-.5) -- 
	                      %(0, \value{grtop}+.5) node[anchor=west,font=#1]{$y$};
	%\foreach \x in {2,4,...,\value{grright}}
		%\draw[tick] (\x,.1) -- (\x,-.1) node[anchor=north,font=#2]{$\x$};
	%\foreach \x in {2,4,...,\value{grleft}}
		%\draw[tick] (-\x,.1) -- (-\x,-.1) node[anchor=north,font=#2]{$-\x\ \ $};
	%\foreach \x in {2,4,...,\value{grtop}}
		%\draw[tick] (.1,\x) -- (-.1,\x) node[anchor=east,font=#2]{$\x$};
	%\foreach \x in {2,4,...,\value{grbot}}
		%\draw[tick] (.1,-\x) -- (-.1,-\x) node[anchor=east,font=#2]{$-\x$};
}

\newcommand{\AxesBlank}[1][\small]{
	\draw[axis] (-\value{grleft}-0.5,0) -- (\value{grright}+0.5,0) node[anchor=south,font=#1]{$x$};
	\draw[axis] (0,-\value{grbot}-0.5) -- (0,\value{grtop}+0.5) node[anchor=west,font=#1]{$y$};
}

\newcommand{\Graph}[3][2]{%
	\draw[step=1,npblue] (-\value{grleft},-\value{grbot}) grid 
											   (\value{grright},\value{grtop});
  \setspace{#1}
	\AxesGen{#2}{#3}
}

\newcommand{\GraphClean}[2]{%
	\draw[step=1,npblue] (-\value{grleft}/\value{divspacex},-\value{grbot}/\value{divspacey}) grid 
											   (\value{grright}/\value{divspacex},\value{grtop}/\value{divspacey});
	\AxesGen{#1}{#2}
}

\newcommand{\AxesOnes}[2]{%
	\setspace{1}
	\AxesGen{#1}{#2}
	%\draw[axis] (-\value{grleft}-1,0) -- 
	                      %( \value{grright}+1,0) node[anchor=south,font=#1]{$x$};
	%\draw[axis] (0,-\value{grbot}-1) -- 
	                      %(0, \value{grtop}+1) node[anchor=west,font=#1]{$y$};
	%\foreach \x in {1,...,\value{grright}}
		%\draw[tick] (\x,.1) -- (\x,-.1) node[anchor=north,font=#2]{$\x$};
	%\foreach \x in {1,...,\value{grleft}}
		%\draw[tick] (-\x,.1) -- (-\x,-.1) node[anchor=north,font=#2]{$-\x\ \ $};
	%\foreach \x in {1,...,\value{grtop}}
		%\draw[tick] (.1,\x) -- (-.1,\x) node[anchor=east,font=#2]{$\x$};
	%\foreach \x in {1,...,\value{grbot}}
		%\draw[tick] (.1,-\x) -- (-.1,-\x) node[anchor=east,font=#2]{$-\x$};
}

\newcommand{\GraphOnes}{%
	\setspace{1}
	\draw[step=1,npblue] (-\value{grleft},-\value{grbot}) grid 
											 (\value{grright},\value{grtop});
  \AxesGen{\small}{\tiny}
}

\newcommand{\GraphHalf}{%
	\draw[step=1,npblue] (-\value{grleft}*2-1,-\value{grbot}*2-1) grid 
											   (\value{grright}*2+1,\value{grtop}*2+1);
	\setspace[2]{1}
  \AxesGen{\small}{\tiny}
}

\newcommand{\GraphSmall}[1][2]{\Graph[#1]{\scriptsize}{\tiny}}
\newcommand{\GraphLarge}[1][2]{\Graph[#1]{\small}{\scriptsize}}
\newcommand{\AxesSmall}{\AxesGen{\scriptsize}{\tiny}}
\newcommand{\AxesLarge}{\AxesGen{\small}{\scriptsize}}

\newcommand{\GraphSmallClean}{\GraphClean{\scriptsize}{\tiny}}
\newcommand{\GraphLargeClean}{\GraphClean{\small}{\scriptsize}}

\newcommand{\ha}[1]{%
  \draw[asymp] (-\value{grleft},#1) -- (\value{grright},#1);
}
\newcommand{\va}[1]{%
  \draw[asymp] (#1,-\value{grtop}) -- (#1,\value{grbot});
}

\newcommand{\GraphIn}[2]{%
	\setgraphsquare{#1}
	\begin{scope}[xshift={(2*\value{cols}-#1-1)*\linespace},yshift={(-#1-1)*\linespace}]
	\AxesGen{\small}{\scriptsize}
	#2
	\end{scope}
}

\newcommand{\GraphInLeft}[2]{%
	\setgraphsquare{#1}
	\begin{scope}[xshift={(#1+1)*\linespace},yshift={(-#1-1)*\linespace}]
	\AxesGen{\small}{\scriptsize}
	#2
	\end{scope}
}
\newcommand{\GraphInOnes}[2]{%
	\setgraphsquare{#1}
	\setspace{1}
	\begin{scope}[xshift={(2*\value{cols}-#1-1)*\linespace},yshift={(-#1-1)*\linespace}]
	\AxesGen{\small}{\scriptsize}
	#2
	\end{scope}
}

\newcommand{\GraphInSet}[1][]{%
	\begin{scope}[xshift={(2*\value{cols}-\value{grright}-1)*\linespace},yshift={(-\value{grtop}-1)*\linespace}]
	\AxesGen{\small}{\scriptsize}
	#1
	\end{scope}
}

\newenvironment{syspic}[2][1.5]{
\begin{icpic}[x=0.5\linespace,y=0.5\linespace]
\useasboundingbox (-.5\textwidth,-5.5) rectangle (0,#2-#1);
\setgraphsquare{#2}
\GraphSmall
}{
\end{icpic}
}

\NewDocumentCommand{\linepoints}{>{\SplitArgument{1}{,}}r() >{\SplitArgument{1}{,}}r()} {
	\linepointsSUB #1 #2
}

\NewDocumentCommand{\linepointslope}{>{\SplitArgument{1}{,}}r() m}{
	\lineslopeSUB #1{#2}
}

\NewDocumentCommand{\lineslopeintercept}{m m}{
	\lineslopeSUB{0}{#2}{#1}
}

\NewDocumentCommand{\linepointsSUB}{m m m m} {%
	\lineslopeSUB{#1}{#2}{(#4-#2)/(#3-#1)}
	%\path[name path=grbound]
		%(-\value{grleft}*\value{countspace}/\value{divspace},
			%-\value{grbot}*\value{countspace}/\value{divspace})
			%rectangle
			%(\value{grright}*\value{countspace}/\value{divspace},
			%\value{grtop}*\value{countspace}/\value{divspace});
	%\path[name path=myline]
		%(-\value{grleft}*\value{countspace}/\value{divspace}-\value{countspace}/\value{divspace},
			%{(#4-#2)/(#3-#1)*(-\value{grleft}-1-#1)*\value{countspace}/\value{divspace}
			%+#2*\value{countspace}/\value{divspace}}) 
			%-- (\value{grright}*\value{countspace}/\value{divspace}+\value{countspace}/\value{divspace},
			%{(#4-#2)/(#3-#1)*(\value{grright}+1-#1)*\value{countspace}/\value{divspace}
			%+#2*\value{countspace}/\value{divspace}});
	%\draw[name intersections = {of = grbound and myline, name = pt},
				%graph] (pt-1) -- (pt-2);
}

\NewDocumentCommand{\lineslopeSUB}{m m m} {%
	\path[name path=grbound]
		(-\value{grleft}*\value{countspacex}/\value{divspacex},
			-\value{grbot}*\value{countspacey}/\value{divspacey})
			rectangle
			(\value{grright}*\value{countspacex}/\value{divspacex},
			\value{grtop}*\value{countspacey}/\value{divspacey});
	\path[name path=myline]
		(-\value{grleft}*\value{countspacex}/\value{divspacex}-\value{countspacex}/\value{divspacex},
			{(#3)*(-\value{grleft}-1-#1)*\value{countspacey}/\value{divspacey}
			+#2*\value{countspacey}/\value{divspacey}}) 
			-- (\value{grright}*\value{countspacex}/\value{divspacex}+\value{countspacex}/\value{divspacex},
			{(#3)*(\value{grright}+1-#1)*\value{countspacey}/\value{divspacey}
			+#2*\value{countspacey}/\value{divspacey}});
	\draw[name intersections = {of = grbound and myline, name = pt},
				graph] (pt-1) -- (pt-2);
}

\NewDocumentCommand{\plotpoint}{>{\SplitArgument{1}{,}}r()}{%
	\plotpointSUB #1
}
\NewDocumentCommand{\plotpointSUB}{m m}{%
	\draw[dotfull] ({(#1)*\value{countspacex}/\value{divspacex}},
		{(#2)*\value{countspacey}/\value{divspacey}})
		circle (\dotsize);
}

\newlength{\dotsize}
\setlength{\dotsize}{0.1\linespace}

%%%%%%%%%%%%%%%%%%%%%%%%%%%%%%%%%%%%%%%%%%%%%%%%%%%%%%%%%%%%%%%%
% Interval Graphing
%%%%%%%%%%%%%%%%%%%%%%%%%%%%%%%%%%%%%%%%%%%%%%%%%%%%%%%%%%%%%%%%
\newcommand{\numline}[3][0]{%
\draw[axis] (#2-1,#1) -- (#3+1,#1);
\foreach \x in {0,1,...,#3}
	\draw[tick] (\x,#1+0.2) -- (\x,#1-0.2) node[anchor=north]{$\x$};
\foreach \x in {#2,...,-1}
	\draw[tick] (\x,#1+0.2) -- (\x,#1-0.2) node[anchor=north]{$\x\phm$};
}

\newcommand{\numlinec}[2][0]{\numline[#1]{-#2}{#2}}

% <---->
\newcommand{\graphaa}[3][0]{%
	\draw[graph] (#2,#1+0.3) -- (#3,#1+0.3);
}

% <----o
\newcommand{\graphao}[3][0]{%
	\draw[graphend,<-] (#2,#1+0.3) -- (#3,#1+0.3);
	\draw[dotempty] (#3,#1+0.3) circle (0.1);
}

% <----*
\newcommand{\graphac}[3][0]{%
	\draw[graphend,<-] (#2,#1+0.3) -- (#3,#1+0.3);
	\draw[dotfull] (#3,#1+0.3) circle (0.1);
}

% o---->
\newcommand{\graphoa}[3][0]{%
	\draw[graphend,->] (#2,#1+0.3) -- (#3,#1+0.3);
	\draw[dotempty] (#2,#1+0.3) circle (0.1);
}

% o----o
\newcommand{\graphoo}[3][0]{%
	\draw[graphend] (#2,#1+0.3) -- (#3,#1+0.3);
	\draw[dotempty] (#2,#1+0.3) circle (0.1);
	\draw[dotempty] (#3,#1+0.3) circle (0.1);
}

% o----*
\newcommand{\graphoc}[3][0]{%
	\draw[graphend,<-] (#2,#1+0.3) -- (#3,#1+0.3);
	\draw[dotempty] (#2,#1+0.3) circle (0.1);
	\draw[dotfull] (#3,#1+0.3) circle (0.1);
}

% *---->
\newcommand{\graphca}[3][0]{%
	\draw[graphend,->] (#2,#1+0.3) -- (#3,#1+0.3);
	\draw[dotfull] (#2,#1+0.3) circle (0.1);
}

% *---->
\newcommand{\graphco}[3][0]{%
	\draw[graphend] (#2,#1+0.3) -- (#3,#1+0.3);
	\draw[dotfull] (#2,#1+0.3) circle (0.1);
	\draw[dotempty] (#3,#1+0.3) circle (0.1);
}

% *----*
\newcommand{\graphcc}[3][0]{%
	\draw[graphend,<-] (#2,#1+0.3) -- (#3,#1+0.3);
	\draw[dotfull] (#2,#1+0.3) circle (0.1);
	\draw[dotfull] (#3,#1+0.3) circle (0.1);
}


\newcommand{\divbracket}[3]{\draw (#1+2*#2+.5,#3) -- (#1+.5,#3) 
                                              arc (20:-20:{.5/sin(20)})}
%%%%%%%%%%%%%%%%%%%%%%%%%%%%%%%%%%%%%%%%%%%%%%%%%%%%%%%%%%%%%%%%
% Division gear
%%%%%%%%%%%%%%%%%%%%%%%%%%%%%%%%%%%%%%%%%%%%%%%%%%%%%%%%%%%%%%%%
%first argument is x-coord of upper left
%second argument is y-coord of upper left
%third argument is number of terms
\newcommand{\synthdiv}[3]{%
  \draw (#1,#2-1) -- ++(1,0) -- ++(0,1); %bracket
	\draw (#1+1.5,#2-2) -- ++(2*#3,0);
	\draw (#1+2*#3-.5,#2-2) -- ++(0,-1)
}
\newcommand{\synthdivtight}[3]{%
  \draw (#1,#2-1) -- ++(1,0) -- ++(0,1); %bracket
	\draw (#1+1.5,#2-2) -- ++(1.5*#3+.5,0);
	\draw (#1+1.5*#3+.25,#2-2) -- ++(0,-1)
}

%%%%%%%%%%%%%%%%%%%%%%%%%%%%%%%%%%%%%%%%%%%%%%%%%%%%%%%%%%%%%%%%
% Commands for 3d drawing
%%%%%%%%%%%%%%%%%%%%%%%%%%%%%%%%%%%%%%%%%%%%%%%%%%%%%%%%%%%%%%%%

\newcommand{\project}[3]{%
	{(#1)*(2+sqrt(2))/4-(#3)*sqrt(5)/2.5},{#2-(#1)*(2-sqrt(2))/4-(#3)*sqrt(5)/5}%
}
\newlength{\dimtd}
\setlength{\dimtd}{2\linespace}

%%%%%%%%%%%%%%%%%%%%%%%%%%%%%%%%%%%%%%%%%%%%%%%%%%%%%%%%%%%%%%%%
% misc math stuff
%%%%%%%%%%%%%%%%%%%%%%%%%%%%%%%%%%%%%%%%%%%%%%%%%%%%%%%%%%%%%%%%
\newcommand{\suchthat}{\;\middle|\;}
\newcommand{\subnew}{_{\mathrm{new}}}
\newcommand{\inv}{^{-1}}


%%%%%%%%%%%%%%%%%%%%%%%%%%%%%%%%%%%%%%%%%%%%%%%%%%%%%%%%%%%%%%%%
% misc other stuff
%%%%%%%%%%%%%%%%%%%%%%%%%%%%%%%%%%%%%%%%%%%%%%%%%%%%%%%%%%%%%%%%
\usepackage{tabularx}
\newcolumntype{C}{>{\centering\arraybackslash}X}
%\newcolumntype{R}{>{\raggedright\arraybackslash}X}
\newcolumntype{R}{>{\hfill\arraybackslash}X}
\usepackage{multirow}

\newcommand{\fitrow}[4][\linespace]{%
	\begin{tikzpicture}[x=#1,y=#1]
		\useasboundingbox ({-.5*(#2)},{-.5*(#3)}) rectangle ({.5*(#2)},{.5*(#3)});
		\draw (0,-1ex) node[anchor=base]{#4};
	\end{tikzpicture}
}

\newlength{\fakeht}
\newcommand{\fakeit}[1]{%
	\setlength{\fakeht}{\totalheightof{#1}}
	\begin{icpic}
	\useasboundingbox (-1in,0) rectangle (1in,\fakeht);
	\draw (0,.5\fakeht) node{#1};
	\end{icpic}
}

\newcommand{\funtable}[4]{%
\[
\renewcommand*{\arraystretch}{1.5}
\begin{array}{>{\color{colf}}l@{\qquad}>{\color{colg}}l@{\qquad}>{\color{colh}}l}
\rpt{#1}{%
#2(\makeblanksmall)=\makeblanksmall&
#3(\makeblanksmall)=\makeblanksmall&
#4(\makeblanksmall)=\makeblanksmall\\}
&
#3(x)=\makeblank[1in]&
#4(x)=\makeblank[1in]
\end{array}
\]
}

\newcommand{\funtablebig}[5]{%
\[
\renewcommand*{\arraystretch}{1.5}
\begin{array}{>{\color{colf}}l@{\quad}>{\color{colg}}l@{\quad}>{\color{colh}}l@{\quad}>{\color{colk}}l}
\rpt{#1}{%
#2(\makeblank[.375in])=\makeblank[.375in]&
#3(\makeblank[.375in])=\makeblank[.375in]&
#4(\makeblank[.375in])=\makeblank[.375in]&
#5(\makeblank[.375in])=\makeblank[.375in]\\}
&
#3(x)=\makeblank[.75in]&
#4(x)=\makeblank[.75in]&
#5(x)=\makeblank[.75in]
\end{array}
\]
}

\newcounter{step}
\newcommand{\steps}[1][1]{\setcounter{step}{#1}}
\newcommand{\numstep}[1][.]{\arabic{step}.\addtocounter{step}{1}}

\newcommand{\functionfamily}[3][]{%
\begin{displaybox}[#2]
\begin{lpic}{}{10}
\GraphIn{4}{#3}
\foreach \y/\lbl in {-1/Form:,
                    -2/Example:,
										-3/Domain:,
										-4/Range:,
										-5/Inverse:,
										-6/Symmetry:,
										-7/Asymptotes:,
										-8/Key Points:,
										-9/Notes:}
	\regtextleft{1}{\y}{\lbl};
\end{lpic}%
\end{displaybox}%
}

\newcommand{\phm}{\phantom{-}}
\newcommand{\ph}{\phantom{1}}
\newcommand{\php}{\phantom{(}}
%%%%%%%%%%%%%%%%%%%%%%%%%%%%%%%%%%%%%%%%%%%%%%%%%%%%%%%%%%%%%%%%
% Fake small caps
%%%%%%%%%%%%%%%%%%%%%%%%%%%%%%%%%%%%%%%%%%%%%%%%%%%%%%%%%%%%%%%%
\usepackage{xstring} % needed for IfEqCase
\usepackage{forloop}

\newcounter{sccounter}
\newcounter{tempStringLength}

\newcommand{\fakesc}[1]{%
    % this \betterfakesc command requires these two packages:
    %   xstring
    %   forloop
    %
    % First, we obtain the length of the input string.
    \StrLen{#1}[\stringLength]%
    %
    % Our main forloop will be using a condition of “while less than \stringLength”,
    % so we’ll need to increase \stringLength by 1 so the forloop will be able to iterate 
    % over the entire string. we’ll use a temporary counter tempStringLength to make
    % this increase. That’s what the next three lines are about.
    \setcounter{tempStringLength}{\stringLength}%
    \addtocounter{tempStringLength}{1}%
    \def\stringLength{\arabic{tempStringLength}}%
    %
    % Here is our main loop. We iterate over the characters in the input string,
    % and the currentLetter is compared to the case rules we have defined. Basically
    % if the currentLetter is any of the lowercase a-z letters, then we apply a 
    % “fake small caps” effect to it and output it.
    \forloop[1]{sccounter}{1}{\value{sccounter}<\stringLength}{%
        \StrChar{#1}{\value{sccounter}}[\currentLetter]%
        %
        \IfEqCase*{\currentLetter}{%
        % The lines below are the rules. Obviously more could be added.
        {a}{{\uppercase{\footnotesize a}}}%
        {b}{{\uppercase{\footnotesize b}}}%
        {c}{{\uppercase{\footnotesize c}}}%
        {d}{{\uppercase{\footnotesize d}}}%
        {e}{{\uppercase{\footnotesize e}}}%
        {f}{{\uppercase{\footnotesize f}}}%
        {g}{{\uppercase{\footnotesize g}}}%
        {h}{{\uppercase{\footnotesize h}}}%
        {i}{{\uppercase{\footnotesize i}}}%
        {j}{{\uppercase{\footnotesize j}}}%
        {k}{{\uppercase{\footnotesize k}}}%
        {l}{{\uppercase{\footnotesize l}}}%
        {m}{{\uppercase{\footnotesize m}}}%
        {n}{{\uppercase{\footnotesize n}}}%
        {o}{{\uppercase{\footnotesize o}}}%
        {p}{{\uppercase{\footnotesize p}}}%
        {q}{{\uppercase{\footnotesize q}}}%
        {r}{{\uppercase{\footnotesize r}}}%
        {s}{{\uppercase{\footnotesize s}}}%
        {t}{{\uppercase{\footnotesize t}}}%
        {u}{{\uppercase{\footnotesize u}}}%
        {v}{{\uppercase{\footnotesize v}}}%
        {w}{{\uppercase{\footnotesize w}}}%
        {x}{{\uppercase{\footnotesize x}}}%
        {y}{{\uppercase{\footnotesize y}}}%
        {z}{{\uppercase{\footnotesize z}}}%
        }%
        % if our \currentLetter isn’t any of the letters we have rules for,
        % then just output it now
        [{\currentLetter}]%
    }%
}

%%%%%%%%%%%%%%%%%%%%%%%%%%%%%%%%%%%%%%%%%%%%%%%%%%%%%%%%%%%%%%%%
% Calculator entry commands
%%%%%%%%%%%%%%%%%%%%%%%%%%%%%%%%%%%%%%%%%%%%%%%%%%%%%%%%%%%%%%%%
\newlength{\inputwidth}
\newlength{\inputheight}

\newcommand{\TIbut}[1]{%
\setlength{\inputwidth}{\widthof{\fakesc{#1}}}%
\setlength{\inputheight}{\heightof{ABC123}}%
\begin{tikzpicture}[baseline=0]%
\useasboundingbox(-3pt,0) rectangle (\inputwidth+3pt,\inputheight);
\draw(0.5\inputwidth,0) node[anchor=base]{\fakesc{#1}};
\draw[rounded corners=3pt] (-3pt,-3pt) rectangle (\inputwidth+3pt,\inputheight+3pt);
\end{tikzpicture}%
}

\newcommand{\TIbutpro}[1]{%
\setlength{\inputwidth}{\widthof{#1}}%
\setlength{\inputheight}{\heightof{ABC123}}%
\begin{tikzpicture}[baseline=0]%
\useasboundingbox(-3pt,0) rectangle (\inputwidth+3pt,\inputheight);
\draw(0.5\inputwidth,0) node[anchor=base]{#1};
\draw[rounded corners=3pt] (-3pt,-3pt) rectangle (\inputwidth+3pt,\inputheight+3pt);
\end{tikzpicture}%
}

\newcommand{\TIright}{%
\setlength{\inputheight}{\heightof{ABC123}}%
\begin{tikzpicture}[baseline=0]%
\useasboundingbox(-0.5\inputheight-3pt,0) rectangle (0.5\inputheight+3pt,\inputheight);
\filldraw (-0.25\inputheight,0) -- (0.25\inputheight,0.5\inputheight) -- (-0.25\inputheight,\inputheight) -- cycle;
\draw[rounded corners=3pt] (-0.5\inputheight-3pt,-3pt) rectangle (0.5\inputheight+3pt,\inputheight+3pt);
\end{tikzpicture}%
}

\newcommand{\TIleft}{%
\setlength{\inputheight}{\heightof{ABC123}}%
\begin{tikzpicture}[baseline=0]%
\useasboundingbox(-0.5\inputheight-3pt,0) rectangle (0.5\inputheight+3pt,\inputheight);
\filldraw (0.25\inputheight,0) -- (-0.25\inputheight,0.5\inputheight) -- (0.25\inputheight,\inputheight) -- cycle;
\draw[rounded corners=3pt] (-0.5\inputheight-3pt,-3pt) rectangle (0.5\inputheight+3pt,\inputheight+3pt);
\end{tikzpicture}%
}

\newcommand{\TIup}{%
\setlength{\inputheight}{\heightof{ABC123}}%
\begin{tikzpicture}[baseline=0]%
\useasboundingbox(-0.5\inputheight-3pt,0) rectangle (0.5\inputheight+3pt,\inputheight);
\filldraw (-0.5\inputheight,0.25\inputheight) -- (0,0.75\inputheight) -- (0.5\inputheight,0.25\inputheight) -- cycle;
\draw[rounded corners=3pt] (-0.5\inputheight-3pt,-3pt) rectangle (0.5\inputheight+3pt,\inputheight+3pt);
\end{tikzpicture}%
}

\newcommand{\TIdown}{%
\setlength{\inputheight}{\heightof{ABC123}}%
\begin{tikzpicture}[baseline=0]%
\useasboundingbox(-0.5\inputheight-3pt,0) rectangle (0.5\inputheight+3pt,\inputheight);
\filldraw (-0.5\inputheight,0.75\inputheight) -- (0,0.25\inputheight) -- (0.5\inputheight,0.75\inputheight) -- cycle;
\draw[rounded corners=3pt] (-0.5\inputheight-3pt,-3pt) rectangle (0.5\inputheight+3pt,\inputheight+3pt);
\end{tikzpicture}%
}

\newcommand{\TI}[1]{\texttt{#1}}

\newcommand{\TIsub}[1]{{\scriptsize #1}}
\newcommand{\TIsup}[1]{{\textsuperscript{#1}}}

\newcommand{\wordor}{\mbox{ or }}
\newcommand{\wordand}{\mbox{ and }}

